\documentclass[11pt,a4paper]{article}
\usepackage[margin=2.5cm]{geometry}
\usepackage{booktabs}
\usepackage{array}
\usepackage{graphicx}

\title{Experimental Materials: Solidarity Game}
\author{Steimanis, Burger, Hayo, Landmann, \& Vollan}
\date{}

\begin{document}
\maketitle

All of you have 200 pesos at the beginning of the game. You will make your decisions on a sheet of paper only, but the decisions that you take are still about real money. For the rest of the game we have formed groups, each consisting of 3 players. Each of the originally invited participants [point to the left side where originally invited participants sit] brought along two friends. One sits in the middle and will play with you [point to the middle]. The other one who is sitting on the right-hand side will not be in the same group. Instead, the third player in your group will be someone from the right-hand side of the room, but you will never exactly know who it is. And the ones on the right-hand side will never know the two other group members they play with. From now on we will call the unknown players ``Player X''.

Whether you can keep the 200-peso given to you or lose them again will depend partly on your choices and partly on your luck. Remember, only one of the games will be randomly selected for the computation of the pay-out at the end of our session. For each group, we now have an opaque bag with 3 balls in it. This means that there are as many balls in the bag as we have players in a group. Each player draws one ball. Out of the 3 balls, there are 2 white balls and 1 red ball. If you draw a white ball you can keep your 200 pesos. If you draw a red ball you lose the 200 pesos you had at the start of the game. This means that one of the three players in each group will lose everything and two out of three will lose nothing.

\begin{center}[Hang up poster with example decision sheet on it]\end{center}

In this game, the two winners can give money to the loser. Before you draw a ball, all of the players will be asked whether and how much they would like to transfer to the other two players in their group in case that they are unlucky, i.e.\ they draw a red ball and lose 200 pesos. Remember that one of the three players will lose for sure. Remember also that there will always be two players in your group who still have their 200 pesos. You can transfer between 0 and 70 of your 200 pesos to the unlucky person in your group. We will ask you to write down on a worksheet how much you would be willing to give to the losing player. Amounts are given in steps of 10 pesos. You can also decide to transfer nothing. Hence, possible transfers are 0, 10, 20, 30, 40, 50, 60 or 70. Every transfer decision you make is as good as any other---there are no wrong decisions. Your transfers will be kept in private, so just choose the amount YOU like best! But remember it is going to be a transfer of real money.

From now on, we will call the group member you know by his or her name (\_\_\_\_\_\_) [ASSISTANTS ENTER THE NAME OF NON-ANONYMOUS PARTNER HERE] and the unknown group member Player X. For the players sitting on the right-hand side [point] there will be two unknown players Player X and Player Y. So, imagine you keep your 200 pesos and Player X loses his 200 pesos. We will ask you to write down on the worksheet how much you would be willing to give to Player X in this case (0, 10, 20, 30, 40, 50, 60 or 70). Now imagine you keep your 200 pesos and the friend you came here with and plays with you in your group loses his or her 200 pesos. Please write down on the worksheet how much you would be willing to give to him or her in this case (0, 10, 20, 30, 40, 50, 60 or 70).

We also want you to think about the transfer of the other winner in your group to the loser. Please guess the amounts that will be transferred. You will earn 10 pesos extra for each correct guess.

Lastly, it is, of course, possible that you draw the red ball and lose. We would like you to guess how much your friend in your group and Player X would be willing to give to you in this case. We will never tell you whether you were right or not. But Lukas will look at the choices made by your friend and Player X and compare their choices to your guess. You will earn 10 pesos extra for each correct guess. The best thing you can do to increase your payoff is to truthfully state what you think y and Player X would do.

\begin{center}[SHOW AND EXPLAIN PARTICIPANT FORM make sure that the player is looking at the form and appears to be sufficiently concentrated]\end{center}

For non-anonymous players:

\begin{center}\includegraphics[width=0.95\textwidth]{figures/transfer_non_anonymous.png}\end{center}

For anonymous Players:

\begin{center}\includegraphics[width=0.95\textwidth]{figures/transfer_anonymous.png}\end{center}

Before we start playing the game, we would like to ask you some questions to see if you understand the game.

\begin{center}
\begin{tabular}{p{11cm}p{4cm}}
\toprule
What is the maximum amount of money you can earn in this game? & [ANSWER: 220] \\
What is the minimum amount of money you can win in this game? & [ANSWER: 0] \\
What is the highest amount you can transfer to the other player? & [ANSWER: 70] \\
What is the least amount of money you can transfer to the losing player? & [ANSWER: 0] \\
How much does the losing player earn, if the other two other players transfer nothing? & [ANSWER: between 0, 10 or 20] \\
\bottomrule
\end{tabular}
\end{center}

Now we will distribute the decision sheets for the second game.

\begin{center}[Distribute decision sheets, collect them and bring them to XX when every participant is finished with filling them out.]\end{center}

\end{document}
